\section{The nonconstant lift on a down, up, down arc in the near range}

% BEGIN K83V1
\paragraph{K83V1. An exact lift in the positive middle letter.}
Use K83J1's notation and guards, and fix the complementary
horizontal target coordinate $k$. Put $i=c-1$, $u=1+p$,
$q=p-1$ and $L_0=\log_p(2h)$. At $i$ let $l$ be the
two-step translated representative of the constant coordinate at
$c-3$ multiplied by $p^\sigma$, obtained by applying one further
Frobenius to U83M6's first-step matrix. This is an element of the
integral two-step Frobenius span. At a primitive $p$th root at $i$
it equals the constant crystalline period $p^\sigma d_i$ through
degree $qM-1$. In fact the propagated modulus is
$\varphi(v^M P_0^{2G})\sim v^M P_0^{qM}P_1^{2G}$.
Here $d_i\in\operatorname{Fil}^{b_1}$: S83M1's up condition
at level one of the final letter, transported back by crystalline
Frobenius, says exactly this at $i$. At $v=0$ the same jet gives
$p^\sigma d_i$ as well.

The target vector
\[
                 x=P_0^{-b_1}P_1^{y}l
\]
is integral at $i$ and has an exact preimage after multiplication
by $p^{E_0}$ in the source invariant Wach lattice at $i$, where
$E_0\le5U$. To see integrality, the lower height $-y$ at $c-2$
gives $P_1^yl\in A\varphi N_{c-2}$. The lower height at $i$
is $-\eta$, so $P_0^{b_1+\eta}$ takes $x$ into $N_i$.
At the primitive root $x$ has a scalar multiple of $d_i$ starting
in degree $-b_1$, plus an error of order at least $qM-b_1$.
The former passes the filtration test; the latter starts above
$\eta$ since $qM>2h$. The integral Taylor criterion gives the
assertion. Its crystalline value at zero is
$p^{\sigma+y-b_1}d_i$.

We give the cost of lifting the \emph{whole} $x$. Set
$H_*=h+\eta$, $c_H=\lceil1+\log_p(H_*+1)\rceil$, and
$a=\lceil H_*p/q^2\rceil+c_H$. All the letter widths are at
most $H_*$. The binomial proof of T83A1 with this $a$ gives
$P_0^a N_i(T_2)\subset F N_i(T_1)$:
for a term $p^{a-j}v^j$ with $qj<H_*$,
$a-j\ge (H_*-qj)p/q^2+c_H$. Here
$a\le H_*p/q^2+3+L_0$ as $H_*+1\le2h$.
The interpolation intervals
\[
 I=[b_1-a+1,b_1],\qquad
 J=[b_1+1,b_1+S+a]\cup[1-\eta,\eta+a]
\]
are disjoint. Indeed $h\le1.1 b_*$, $H_*\le1.15b_*$ and
$\eta\le b_*/20$ on the given mixed word.
Since $p/q^2\le11/100$ and $6+2L_0\le6U\le6b_*/160$,
we have $2a+\eta<b_*\le b_1$.

For $D=(\gamma-1)/v$ the polynomial $P_I(D)$ of T83A4 sends $x$
to $P_0^aN_i$. The same triangular calculation works here:
all terms below $qM-b_1$ are along $d_i$, the $a$ factors remove
the first $a$ pole degrees, and $qM-b_1-a\ge\eta$ handles the
other terms after division by $P_0^a$. Lift $P_I(D)x$ and
$P_0^a x$. Applying $P_J(D)$ to the second lift divided by
$P_0^a$ gives an integral lift of $P_J(D)x$ by the source
pole-jet calculation in T83A4, with the central jumps now
$[-\eta,\eta]$. As before, interpolation gives
$n_1\in N_i(T_1)$ with $Fn_1=p^{E_0}x$ at
\[
 E_0=\left\lfloor(3a+S+2\eta-1)/q+L_0\right\rfloor.
\]
Every difference in the remote interval count is at most
$b_1+\eta\le2h$, so the indicated logarithm suffices. For the
numerical bound, $3p^3/q^3\le4$, $p/q\le1.1$ and $p\ge11$ give
\[
 E_0\le 4h/p^2+4\eta/p^2+(S+2\eta)/q+
                     (9+3L_0)/q+L_0
       \le 5\{h/p^2+(w+S+\eta)/p+L_0+1\}=5U.
\]
This estimate precedes the final cut lift and is independent of
the particular scalar $f$.
% END K83V1

% BEGIN K83V2
\paragraph{K83V2. Cancelling the primitive constant and lifting the remaining scalar.}
There is a scalar $t_3\in O$ with finite $v_p(t_3)\le14U$ such
that $t_3f_2$ of K83J1 satisfies the actual $\mathrm{LiftsAt}$
at exponent $R+\sigma\le6U$. To construct it, write
$f_2=P_2^y g/P_1^{b_1}$ with $g\in A$. The target vector
\[
              V=f_2\varphi l=g\varphi x
\]
lies in the one-step translated span $A\varphi N_i(T_2)$ since
$x$ is integral by K83V1. It represents $p^\sigma f_2$ in
coordinate $k$: $\varphi l$ is precisely U83M6's three-step
forward column, and $\Lambda f_2\in A$. In particular this
representation has the required integral numerator.

Use the one-step horizontal comparison T83A1 on the source basis
at $i$, acting on the specialization of $n_1$. With modulus
$vP_0^B$ and exponent $s_1=\lceil(1+B)/q+1\rceil$, this gives
$Y\in A\varphi N_i(T_1)$ with primitive Taylor value
$p^{s_1}\varphi\operatorname{cmp}(n_1)$ modulo $t^B$. Its
gamma difference is in $vP_0^B A\varphi N_i(T_1)$.
Indeed it vanishes through order $B-1$ at each primitive branch,
since the indicated crystalline period is fixed, and its
translated-basis coordinates vanish at $v=0$ since gamma acts
trivially on the constant fibre. The translated basis is regular
and invertible in the primitive Taylor coordinates, by composition
with Frobenius at zero. Therefore distinguished division yields
the product $vP_0^B$ integrally. The lower height at $c$ implies
that $P_0^B Y$ is integral and $(\gamma-1)Y\in vN_c(T_1)$.

Write $d_c=\varphi d_i$. At the primitive $p$th root of the
final letter, the vectors $V$ and $FY$ are respectively congruent
through degree $B-1$ to the constants
\[
     p^{\sigma-\kappa}X_0d_c,\qquad
     p^{s_1+E_0+\sigma+y-b_1}d_c .
\]
For the first equality, $P_1$ is a Taylor unit and K83J1 gives
$f_2\equiv X_0/p^\kappa$ there. The column $\varphi l$ is
constant through degree $qM-1\ge B-1$ by its propagated modulus.
The second equality is K83V1's specialization and the source
truncation. Put $D_1=s_1+E_0+y-b_1+\kappa$. If $X_0=0$, take
$t_3=1$ and $b'=0$. Otherwise choose $t_3,b'\in O$ such that
$t_3X_0=p^{D_1}b'$: take $t_3=1$, $b'=p^{-D_1}X_0$ if the
latter is integral, and $t_3=p^{D_1}/X_0$, $b'=1$ otherwise.
Then $v_p(t_3)=\max(0,D_1-v_p(X_0))$. In the nontrivial case
K83J2 bounds this by
\[
 E_0+s_1-B/p+B/p^2+2+L_0
 \le 5U+2.1h/p^2+4.1+L_0\le14U .
\]
Here $B\le h$, $p/q\le1.1$, and
$s_1<(B+1)/q+2$.

The vector $D'=t_3V-b'FY$ is integral in the one-step translated
basis and its primitive Taylor coefficients below $B$ vanish
by construction. Distinguished division in that same regular
basis gives $D'\in P_0^B A\varphi N_i(T_2)\subset N_c(T_2)$.
At this cut lift $p^RD'$ to an invariant $z'\in N_c(T_1)$
by U83M1. The source vector $z'+p^R b'Y$ now maps exactly to
$p^R t_3 V$. It has finite poles and gamma difference in
$vN_c(T_1)$. Its image has the required translated numerator and
represents $p^{R+\sigma}t_3f_2$ in coordinate $k$, with zero
physical remainder. This proves the claimed native scalar lift.

Finally K83J3 lifts the separated integral constant at the same
exponent. Add that lift times $t_3$ and use
$p^Kf=X_2+f_2$. The scalar $t_3p^K$ divides
$p^{\lceil K+v_p(t_3)\rceil}$ in $O$, so the original $f$
lifts at an integer exponent at most
\[
 6U+14U+1+(\eta+2+L_0)\mathbf1_{\rm band}
       \le22U+1+\eta\mathbf1_{\rm band}.
\]
This covers every complementary coordinate on the actual
down, up, down arc in the guard $s_{c-1}<p(s_{c-2}+S+\eta)$,
for every tested surjection and coefficient extension.
% END K83V2
